% Containers, objects, calls, and traps. Requires tikz_styles.tex.
\begin{figure}[p]
\centering
\begin{tikzpicture}[node distance=1.5mm and 3.5mm]
  % LIST
  \node[font=\tiny\sffamily\bfseries, anchor=west] at (0,0) {LIST};
  \node[slot, fill=pyhi, minimum width=34mm, anchor=north west] (lh) at (0,-0.25)
    {capacity $|$ length};
  \node[slot, fill=pyhi, minimum width=34mm, below=1mm of lh] (lb)
    {0 $|$ ob\_item};
  \node[slot, fill=pymem, minimum width=34mm, below=2.2mm of lb] (lv)
    {element value};
  \node[slot, fill=pymem, minimum width=34mm, below=1mm of lv] (lt)
    {element tag};
  % TUPLE
  \node[font=\tiny\sffamily\bfseries, anchor=west] at (4.6,0) {TUPLE};
  \node[slot, fill=pyrf, minimum width=34mm, anchor=north west] (th) at (4.6,-0.25)
    {handle: size $|$ addr};
  \node[slot, fill=pymem, minimum width=34mm, below=2.2mm of th] (tv)
    {element value};
  \node[slot, fill=pymem, minimum width=34mm, below=1mm of tv] (tt)
    {element tag};
  % DICT
  \node[font=\tiny\sffamily\bfseries, anchor=west] at (9.2,0) {DICT};
  \node[slot, fill=pyink, minimum width=42mm, anchor=north west] (dh) at (9.2,-0.25)
    {slot\_count $|$ used};
  \node[slot, fill=pyink, minimum width=42mm, below=1mm of dh] (dv)
    {version $|$ order\_len};
  \node[slot, fill=pyink, minimum width=42mm, below=1mm of dv] (dp)
    {order\_ptr $|$ table\_ptr};
  \node[slot, fill=pymem, minimum width=42mm, below=2mm of dp] (do)
    {order: key value, key tag};
  \node[slot, fill=pyacc, minimum width=42mm, below=1mm of do] (dt)
    {table slot 64\,B: key, val};
\end{tikzpicture}
\caption{Sequence and dict layout, 16-byte slots.
A list is a stable 32-byte object plus a relocatable element buffer.
Growing the buffer does not move the handle.
Each element is a value slot and a tag slot (32 bytes).
An empty list allocates no buffer.
A tuple stores its length in the handle, so it has no header.
A dict is a stable 48-byte object, an insertion-order key sidecar, and a relocatable open-addressed table.
Empty and deleted table slots are \texttt{UNINIT} and \texttt{TOMBSTONE}.
New keys append to the sidecar and bump \texttt{version}; a value overwrite does neither.
Iteration of a dict follows the sidecar, never hash-slot order.}
\label{fig:layouts}
\end{figure}

\begin{figure}[htbp]
\centering
\begin{tikzpicture}[node distance=1.4mm]
  \node[slot, fill=pyhi, minimum width=86mm] (h)
    {obj$+0$ \quad ob\_kind \quad ob\_flags \quad ob\_type};
  \node[slot, fill=pyres, minimum width=86mm, below=of h] (t)
    {obj$+16$ \quad self tag \quad PY\_TAG\_OBJECT};
  \node[slot, fill=pymem, minimum width=86mm, below=of t] (f)
    {obj$+32$ onward \quad field value, field tag, stride 32\,B};
\end{tikzpicture}
\caption{General heap object.
The handle is \texttt{OBJECT} plus the object address.
The kind is not in the tag: \texttt{LOAD\_ATTR} and \texttt{CALL} read \texttt{ob\_head} first.
Kinds in use are instance (dict at field 0), type (dict, base, name), bound method, builtin, exception, cell, and function (code object plus a cell tuple).
Lists, dicts, and sets are \texttt{MUT\_COLLEC}, not \texttt{OBJECT}.
Strings are \texttt{SHORT\_STR} or \texttt{LONG\_STR}.
Range has its own tag.}
\label{fig:object}
\end{figure}

\begin{figure}[htbp]
\centering
\begin{tikzpicture}[node distance=2.4mm and 3mm]
  \node[ctl] (op) {opcode};
  \node[acc, right=of op] (init) {CP\_INIT};
  \node[mem, right=of init] (hdr) {CP\_HDR};
  \node[mem, below=of hdr] (val) {CP\_VAL};
  \node[mem, left=of val] (tag) {CP\_TAG};
  \node[ctl, left=of tag] (done) {CP\_DONE};
  \node[acc, below=5mm of tag] (probe) {CP\_DICT\_PROBE};
  \node[ink, right=of probe] (attr) {CP\_ATTR\_*};
  \draw[arr] (op) -- (init);
  \draw[arr] (init) -- (hdr);
  \draw[arr] (hdr) -- (val);
  \draw[arr] (val) -- (tag);
  \draw[arr] (tag) -- (done);
  \draw[arr] (hdr.south) -- ++(0,-0.15) -| (probe.north);
  \draw[arr] (probe) -- (attr);
  \draw[arr] (done.north) -- ++(0,0.45) -| (op.north);
\end{tikzpicture}
\caption{Container microcode (\texttt{S\_CONTAINER}).
Decode sets \texttt{dec\_is\_container} and execute transfers here, skipping \texttt{S\_MEM} and \texttt{S\_WB}.
The phase register is \texttt{container\_phase\_r}.
Shared phases read or write a header, a value slot, and a tag slot, then \texttt{CP\_DONE} returns to fetch.
Dict and set operations insert a probe loop.
Attribute operations insert a walk of the instance dict and the type MRO.
The arms live in the \texttt{pycore\_cont\_*.svh} files included into the core.
Fifty-six of the sixty-four \texttt{CONT\_*} codes are in use.
A recoverable failure sets \texttt{trap\_marshal\_pending\_r} and \texttt{CP\_DONE} leaves for the mailbox instead of fetch.}
\label{fig:contfsm}
\end{figure}

\begin{figure}[p]
\centering
\begin{tikzpicture}[node distance=2mm and 2.4mm]
  \node[ctl, minimum width=28mm, font=\tiny\sffamily] (py) {PyCore};
  \node[ex, minimum width=28mm, font=\tiny\sffamily, right=14mm of py] (ex) {excore};
  \node[mem, font=\tiny\sffamily, below=of py, minimum width=34mm, align=left] (pyl)
    {hash and probe\\rich equality\\tombstone skip\\spare-cap append\\last-element delete\\empty extend\\tuple scan};
  \node[ex, font=\tiny\sffamily, below=of ex, minimum width=36mm, align=left] (exl)
    {LIST\_GROW append\\LIST\_EXTEND\\mid-list delete\\DICT\_GROW rehash\\SET\_GROW\\uncontaminated\\SET/DICT update};
  \draw[arr] (py) -- (pyl);
  \draw[arr] (ex) -- (exl);
  \draw[darr] (pyl.east) -- node[anno, above, font=\tiny] {before any commit} (exl.west);
\end{tikzpicture}
\caption{What stays on PyCore, and what is finished by excore.
Hashing, equality, and the open-addressed probe stay on the hart.
Capacity changes and the long memmoves go to firmware, which must complete the semantic effect rather than resize and retry.
The trap is raised before any register-file, heap, or data-memory commit, which is what makes a later \texttt{RETRY} legal.
No current handler returns \texttt{RETRY}; they return \texttt{COMPLETED}.
Contaminated containers (an \texttt{OBJECT} key or element) and tuple-sourced set updates stay on PyCore in \texttt{pycore\_cont\_bulk}.}
\label{fig:split}
\end{figure}

\begin{figure}[htbp]
\centering
\begin{tikzpicture}[node distance=2.6mm and 3mm]
  \node[ctl] (call) {CALL / CALL\_KW\\CALL\_FUNCTION\_EX};
  \node[mem, right=of call] (codc) {CODC or\\dmem fields};
  \node[acc, right=of codc] (bind) {shared binder\\args, defaults,\\*args, **kwargs};
  \node[rf, below=of bind] (spill) {spill if needed};
  \node[ink, left=of spill] (frame) {push 32\,B frame};
  \node[ctl, left=of frame] (go) {redirect fetch\\to entry slot};
  \draw[arr] (call) -- (codc);
  \draw[arr] (codc) -- (bind);
  \draw[arr] (bind) -- (spill);
  \draw[arr] (spill) -- (frame);
  \draw[arr] (frame) -- (go);
\end{tikzpicture}
\caption{Call, as it sits among the other machines.
The phase-by-phase binder is the sibling note \texttt{call\_fsm.tex}.
A user function is a code object, except a closure, which is an \texttt{OBK\_FUNCTION} wrapping the code object and a cell tuple.
\texttt{CALL} checks the callable, binds positional and keyword arguments against \texttt{co\_varnames}, then enters the frame manager.
Builtins with a hardware fast path (\texttt{len}, \texttt{range}, \texttt{set}, two-argument \texttt{max}, \texttt{int} and \texttt{str} conversion) dispatch inside this FSM.
Other builtin calls raise the recoverable \texttt{PY\_TRAP\_BUILTIN\_CALL}.
Return pops the descriptor, writes the result into the caller's ring, and fills any spilled prefix the caller needs.}
\label{fig:call}
\end{figure}

\begin{figure}[htbp]
\centering
\begin{tikzpicture}[node distance=2.4mm]
  \node[ink, minimum width=72mm] (n0)
    {node$+0$ \quad previous pointer};
  \node[ink, minimum width=72mm, below=of n0] (n1)
    {node$+16$ \quad saved active exception};
  \node[ctl, below=4mm of n1] (ops)
    {PUSH\_EXC\_INFO \quad CHECK\_EXC\_MATCH \quad POP\_EXCEPT \quad RAISE};
\end{tikzpicture}
\caption{Exception-info stack.
The arena is 4\,KB at \texttt{0xF00000}, 32 bytes per node, depth 128.
It is a data-memory LIFO owned by \texttt{pycore\_exc\_stack}, separate from both the operand ring and the call-frame stack.
The native-method table and the \texttt{StopIteration} sidecar sit in the same arena.
\texttt{FOR\_ITER} exhaustion compares the raised type by identity with the sidecar latched at boot, and redirects to \texttt{POP\_ITER} without a user handler.
A user \texttt{raise} with no handler is still fatal \texttt{PY\_TRAP\_RAISE}.}
\label{fig:exc}
\end{figure}

\begin{figure}[htbp]
\centering
\begin{tikzpicture}[node distance=2.2mm and 4mm]
  \node[trap, minimum width=38mm, align=left, font=\tiny\sffamily] (fat)
    {\textbf{fatal}\\
     type, stack, div0\\
     FPU, illegal, call filter\\
     mem fault, align\\
     attr error, raise\\
     overflow, value};
  \node[ex, minimum width=42mm, align=left, font=\tiny\sffamily, right=of fat] (rec)
    {\textbf{recoverable}\\
     list grow, extend, delete\\
     dict grow, update, merge\\
     set grow, set update\\
     builtin call, slice};
  \node[ctl, minimum width=28mm, below=4mm of fat] (halt) {S\_HALT};
  \node[ex, minimum width=36mm, below=4mm of rec] (mb) {mailbox};
  \draw[arr] (fat) -- (halt);
  \draw[arr] (rec) -- (mb);
\end{tikzpicture}
\caption{Trap classes.
\texttt{pycore\_trap\_recoverable} in \texttt{pycore\_defs.svh} is the split.
With excore disabled, a recoverable code is fatal too.
\texttt{COMPLETED} applies the popped and pushed entries and resumes at the next instruction.
\texttt{RETRY} redirects fetch to the same PC.
\texttt{FATAL} enters \texttt{pycore\_trap} with the firmware's code.
\texttt{NEED\_HEAP} collects if the collector is enabled, and halts otherwise.}
\label{fig:traps}
\end{figure}

\begin{figure}[htbp]
\centering
\begin{tikzpicture}[node distance=1.3mm]
  \node[slot, fill=pyhi, minimum width=88mm] (b0) {\texttt{0x3E0} \quad module code object};
  \node[slot, fill=pyhi, minimum width=88mm, below=of b0] (b1) {\texttt{0x400} \quad globals dict};
  \node[slot, fill=pyhi, minimum width=88mm, below=of b1] (b2) {\texttt{0x420} \quad builtins dict};
  \node[anno, below=1mm of b2] {each pair is a 16-byte value slot and a 16-byte tag slot};
\end{tikzpicture}
\caption{Boot record, read by \texttt{S\_BOOT} when \texttt{BOOT\_EN=1}.
The walker checks the tags, caches \texttt{co\_consts} and \texttt{co\_names}, latches the globals and builtins bases, and redirects fetch to the module entry slot.
A serialized code object is eight tagged fields: entry slot, consts, names, packed metadata, defaults, varnames, keyword defaults, and the raw exception table.
\texttt{LOAD\_CONST} indexes \texttt{co\_consts} with two data reads (value, then tag).
\texttt{LOAD\_GLOBAL} and \texttt{LOAD\_NAME} probe the frame globals and then the builtins dict; a miss is \texttt{MEM\_FAULT}.}
\label{fig:boot}
\end{figure}
